% =============================================================
%  ch12_simulationswerkzeuge.tex  --  Simulationswerkzeuge im ROS 2-Ökosystem
% =============================================================

\chapter{Simulationswerkzeuge im ROS\,2-Ökosystem}

\textbf{Gazebo Classic ist seit Januar 2025 EOL}; ``modernes Gazebo''
(früher Ignition) ist der Nachfolger:
\begin{longtable}{@{}p{3.2cm}p{4.4cm}p{7.4cm}@{}}
\toprule
\textbf{Werkzeug} & \textbf{Charakter} & \textbf{Eignung} \\
\midrule
\textbf{Gazebo Harmonic} & Standard-Simulator für Jazzy, Bridge \code{ros\_gz} & Allrounder: Physik, Sensoren (Kamera, Tiefe, IMU, LiDAR), \code{gz\_ros2\_control}. \textbf{Erste Wahl.} \\
NVIDIA Isaac Sim         & photorealistisch, GPU & Vision/Pose-Schätzung mit realistischen Bildern; braucht starke GPU; nur LTS-Distros \\
Webots                   & einfach, didaktisch    & schneller Einstieg, gute ROS\,2-Anbindung \\
MuJoCo                   & schnelle, präzise Dynamik & kontaktreiche Regelung, RL; weniger Sensor-Realismus \\
PyBullet                 & leichtgewichtig        & schnelles Prototyping von Kinematik/Dynamik \\
\bottomrule
\end{longtable}
\noindent \textbf{RViz2} ist kein Simulator, sondern Visualisierung -- du nutzt
es parallel.

\section{Was lässt sich vorab simulieren?}
Sehr viel -- die meiste \emph{Logik} der vier Projekte ist simulierbar:
\begin{itemize}
\item \textbf{Kinematik/Dynamik} des Stepper-Roboters: Gelenkbewegung,
  Trajektorienabfahrt, Kollisionen, Massen/Trägheiten.
\item \textbf{Sensoren:} simulierte Stereokamera (gerenderte Bilder -- du kannst
  deine Merkmalserkennung darauf testen), Tiefenbild, IMU, LiDAR, Odometrie.
\item \textbf{Tischtennisball:} Projektilflug inkl.\ Gravitation und Aufprall
  als Starrkörperphysik; dein Kalman-Filter kann auf den simulierten
  3D-Positionen laufen.
\item \textbf{Navigation:} Nav2 + simulierte Odometrie/LiDAR in einer Welt.
\item \textbf{Greifen:} MoveIt\,2-Planung in simulierter Szene mit platzierten
  Bauteilen.
\end{itemize}

\begin{infobox}[title=Kannst du die Projekte vollständig in Simulation umsetzen?]
\textbf{Die Steuerungs-, Planungs- und Wahrnehmungs-Logik: ja.} Du kannst
Trajektorienregelung, Kalman-Schätzung, Navigation und die Vision-Pipeline
komplett in Gazebo entwickeln und testen. \textbf{Was die Simulation nicht
abbildet:} reales Kamerarauschen/Kalibrierfehler, Schrittverluste und
Resonanzen des realen Steppers, exakte Kontaktdynamik (Ballspin/-absprung sind
Näherungen), reale Latenzen und Zeitverhalten. Diese Effekte treten erst auf
der echten Hardware zutage.
\end{infobox}

% =============================================================
\section{MuJoCo und Isaac Sim im Detail}
Zwei Werkzeuge der obigen Tabelle verdienen eine genauere Betrachtung, weil sie
in der modernen Robotik -- besonders im Lernen von Regelungen
(\emph{Reinforcement Learning}, RL) und in der Wahrnehmung -- eine
Sonderstellung einnehmen. Sie lösen Probleme, an denen Gazebo als
Allzweck-Simulator an seine Grenzen stößt: \textbf{MuJoCo} bei der schnellen,
präzisen Kontaktdynamik, \textbf{Isaac Sim} beim photorealistischen Rendern und
beim massiv parallelen GPU-Training.

\subsection{MuJoCo -- Multi-Joint dynamics with Contact}
\textbf{MuJoCo} (\emph{Multi-Joint dynamics with Contact}) ist eine
quelloffene Physik-Engine, ursprünglich 2012 von Emanuel Todorov vorgestellt
\cite{todorov2012mujoco} und seit 2021 von Google DeepMind entwickelt und unter
Apache-2.0 frei verfügbar \cite{mujoco_docs}. Ihr Markenzeichen ist ein
\textbf{stabiler, weicher Kontaktlöser} (konvexe Optimierung pro Zeitschritt),
der kontaktreiche Vorgänge -- Greifen, Laufen, Anschlagen -- numerisch stabil
und schnell rechnet.

\begin{itemize}
\item \textbf{Stärken:} sehr schnelle und genaue Dynamik (oft schneller als
  Echtzeit, hunderte bis tausende Schritte/s), differenzierbare Physik,
  analytisch saubere Gelenk-/Sehnenmodelle, exzellent für Regelungs- und
  RL-Forschung. Mit \textbf{MJX} läuft MuJoCo über JAX auf der GPU und
  simuliert tausende Welten parallel.
\item \textbf{Schwächen:} \emph{kein} photorealistisches Rendering (einfacher
  OpenGL-Viewer), weniger fertige Sensor-Plugins als Gazebo, Modelle in eigenem
  \code{MJCF}-XML (URDF-Import möglich, aber oft mit Nacharbeit).
\item \textbf{Robotik-Relevanz:} De-facto-Standard im RL-Benchmarking
  (\code{Gymnasium}/\code{dm\_control}), Lokomotion vierbeiniger Roboter,
  geschickte Manipulation, Modellprädiktive Regelung (MPC). Wenn die
  \emph{Kontaktphysik} und die \emph{Trainingsgeschwindigkeit} zählen, ist
  MuJoCo erste Wahl.
\end{itemize}

\subsection{NVIDIA Isaac Sim -- photorealistisch und GPU-parallel}
\textbf{Isaac Sim} ist NVIDIAs Robotik-Simulator auf Basis der
\textbf{Omniverse}-Plattform und des \code{USD}-Szenenformats
\cite{isaacsim_docs}. Die Physik liefert \textbf{PhysX\,5} (GPU-beschleunigt),
das Rendering nutzt \textbf{RTX-Raytracing} -- also physikalisch plausibles
Licht, Schatten und Materialien.

\begin{itemize}
\item \textbf{Stärken:} \emph{photorealistische} Kamerabilder, exzellente
  Erzeugung synthetischer Trainingsdaten (\emph{Synthetic Data}, inkl.\
  automatischer Labels: Segmentierung, Bounding-Boxes, Tiefe), \emph{Domain
  Randomization} für robuste Vision-Modelle, GPU-parallele Physik für
  großskaliges RL über \textbf{Isaac Lab} \cite{mittal2023isaaclab}.
\item \textbf{Schwächen:} braucht eine starke NVIDIA-RTX-GPU (VRAM-hungrig),
  hohe Einstiegshürde, große Installation; weniger ``leichtgewichtig'' als
  Gazebo oder MuJoCo für schnelle Iteration.
\item \textbf{Robotik-Relevanz:} überall dort, wo das \emph{Aussehen} der Szene
  entscheidend ist -- Training von Objekterkennung/Pose-Schätzung, Sim-to-Real
  für kamerabasierte Politiken, digitale Zwillinge ganzer Fertigungszellen.
\end{itemize}

% =============================================================
\section{Lassen sie sich integriert einsetzen?}
\textbf{Ja -- aber meist arbeitsteilig, nicht als ein gemeinsamer
Simulationskern.} Beide sprechen ROS\,2, sodass sie sich in dieselbe Knoten- und
Topic-Architektur einfügen wie Gazebo. Der gemeinsame Nenner ist wieder
\code{ros2\_control} (vgl.\ das folgende Kapitel zur Sim-Real-Trennung): die
High-Level-Logik bleibt identisch, nur die unterste Schicht -- der Simulator
hinter dem Hardware-Interface -- wird getauscht.

\begin{longtable}{@{}p{3.0cm}p{5.4cm}p{6.6cm}@{}}
\toprule
\textbf{Werkzeug} & \textbf{ROS\,2-Anbindung} & \textbf{Typische Rolle im Workflow} \\
\midrule
Gazebo Harmonic & \code{ros\_gz}-Bridge, \code{gz\_ros2\_control} & Allround-Integrationstest, Sensorik, Nav2/MoveIt2 \\
MuJoCo          & \code{mujoco\_ros2\_control}, eigene Bridges    & Regler-/RL-Training, Kontaktdynamik, MPC \\
Isaac Sim       & natives \code{ROS\,2-Bridge}-Extension         & Vision-Trainingsdaten, photorealistische Pose-Schätzung, RL via Isaac Lab \\
\bottomrule
\end{longtable}

\noindent Der heute verbreitete \textbf{kombinierte Workflow} sieht so aus:
\begin{enumerate}
\item \textbf{Politik/Regler lernen} -- massiv parallel in \textbf{Isaac Lab}
  (auf Isaac Sim) oder in \textbf{MuJoCo/MJX} auf der GPU. Hier zählen
  Geschwindigkeit und Kontaktphysik, nicht das Aussehen.
\item \textbf{Wahrnehmung trainieren} -- photorealistische, automatisch
  gelabelte Bilder aus \textbf{Isaac Sim} für die Objekt-/Pose-Erkennung
  (Domain Randomization gegen den Sim-to-Real-Gap).
\item \textbf{Integrieren und gegentesten} -- die fertigen Knoten im
  \textbf{Gazebo}-Allroundmodell oder direkt am realen Roboter über
  \code{ros2\_control} verifizieren.
\end{enumerate}

\begin{infobox}[title=Einordnung für dieses Projekt]
Für den Stepper- und Tischtennisroboter ist \textbf{Gazebo Harmonic die
praktische Basis}: Kinematik, Sensorik, Kontakt Ball/Schläger, ros2\_control.
\textbf{MuJoCo} lohnt sich, falls die
\emph{Schlag-/Kontaktdynamik} oder eine \emph{gelernte Schlagpolitik}
(RL) genauer untersucht werden soll -- es rechnet kontaktreiche Vorgänge
schneller und stabiler. \textbf{Isaac Sim} wird relevant, wenn die
\emph{Ballerkennung} mit photorealistischen, gelabelten Bildern trainiert
werden soll, statt sie nur an realen Kamerabildern abzustimmen. Beide setzen
voraus: MuJoCo profitiert von einer GPU (MJX), Isaac Sim \emph{erfordert} eine
starke NVIDIA-RTX-GPU -- das ist die wichtigste Hürde gegenüber dem leichten
Gazebo-Setup.
\end{infobox}
